% TOP_PROBLEM: {"id": 3398, "title": "Subset-sums equality (pigeonhole version)", "category_id": 15, "category": {"id": 15, "name": "computer_science", "display_name": "Computer Science", "description": "Computational complexity, algorithms, and theoretical CS.", "slug": "computer-science", "order_index": 15, "created_at": "2026-07-31T15:26:25.671Z"}, "status": "open", "problem_number": "OPG-163", "set_id": 12}
% FIRST_POSTED: 2026-10-02
% ATTEMPT_STATUS: unresolved
% UnsolvedMath ID 3398; source catalogue code OPG-163
% !TeX program = lualatex
% GPT-6 Astra Ultra research attempt; independently checked partial work, 2026-10-02.
\documentclass[11pt,a4paper]{article}
\usepackage[margin=25mm]{geometry}
\usepackage{fontspec}
\setmainfont{lmroman10-regular.otf}[BoldFont=lmroman10-bold.otf,
 ItalicFont=lmroman10-italic.otf,BoldItalicFont=lmroman10-bolditalic.otf]
\usepackage{amsmath,amssymb,mathtools}
\usepackage{unicode-math}
\setmathfont{latinmodern-math.otf}
\AtBeginDocument{\let\setminus\smallsetminus}
\usepackage{xurl,hyperref}
\hypersetup{colorlinks=true,urlcolor=blue}
\setcounter{secnumdepth}{0}
\setlength{\parindent}{0pt}
\setlength{\parskip}{5pt}
\setlength{\emergencystretch}{3em}
\begin{document}
\section*{Subset-sums equality (pigeonhole version)}\label{3398}
{\small UnsolvedMath ID 3398 · OPG-163 · Computer Science · OpenGarden}
\subsection{Definitions and mathematical statement}
The input consists of positive integers $a_1,\ldots,a_n$ written in
binary, with the promise
$S=\sum_i a_i<2^n-1$. For each index set $I\subseteq[n]$, its subset
sum is $s(I)=\sum_{i\in I}a_i$, with $s(\varnothing)=0$.
There are $2^n$ subsets but only $S+1<2^n$ possible integer sums,
so two distinct subsets necessarily have the same sum.

The search problem is to output distinct index sets $I,J$ such that
\[
                         \sum_{i\in I}a_i=\sum_{j\in J}a_j.
\]
Equivalently one may require both sets to be nonempty and disjoint:
remove their common indices, and positivity ensures that neither
remaining side is empty. Is there one deterministic algorithm that
always finds such a pair on promised inputs in time polynomial in
the total input bit length?

The promise implies each $a_i<2^n$, so total bit length is at most
quadratic in $n$; polynomial time in either size convention is
equivalent on the promised domain. Equal values at different indices
are permitted. Merely deciding existence gives no search solution,
since the promise already makes that answer yes. Approximate
equality or a running time polynomial in $S$ does not meet the
exact polynomial-bit-time target.
\subsection{Short English statement}
Can equal sums from two different subsets be found in polynomial time whenever a simple counting bound guarantees their existence?
\subsection{Research attempt: exact collision dynamic programming and bit complexity}
\paragraph{Scope and literature check.}
GPT-6 Astra Ultra, 2 October 2026. Values are positive binary integers,
indices distinguish equal input values, and the goal is a deterministic
algorithm polynomial in input bit length. Jin and Wu [S3] give faster
exponential algorithms for this pigeonhole equal-sums search problem.
The author publication list [S4] records a 2026 follow-up with
subexponential and lower-bound results; that listing alone is not a
proof of a polynomial algorithm for the full promise here. The work
below gives an exact deterministic search procedure and carefully
isolates the numeric-size dependence that prevents it from answering
the general question.

\paragraph{1. Detecting the first collision, with witnesses.}
Maintain a table mapping each attainable sum to one subset of the
indices processed so far, represented by a bit mask. Initially the
only entry is $0\mapsto\varnothing$. Before processing $a_i$, take
a snapshot of the current table. For each old entry $s\mapsto I$,
the new subset $I\cup\{i\}$ has sum $t=s+a_i$. If $t$ is already
an old attainable sum with witness $J$, output these two subsets.
They are different because one contains $i$ and the other does not.
If there is no such intersection, retain all old entries and add
the shifted entries. Two shifted entries cannot collide with each
other, since addition of the same integer preserves distinct sums.

Inductively, before the first collision every subset of the processed
prefix has a distinct table entry and the table size doubles at
each step. If the procedure finished all $n$ steps without collision,
it would contain $2^n$ different integers between zero and $S$,
contrary to the promise $S+1<2^n$. Thus the procedure always returns
an exact equal-sum pair on promised inputs. Removing their common
indices gives disjoint witnesses of equal sum; positivity implies
both remaining sides are nonempty.

\paragraph{2. What running time is actually proved.}
All table keys lie in $[0,S]$. A direct array implementation with
snapshots takes $O(n(S+1))$ arithmetic and table operations. Keys
have $O(\log(S+1))$ bits and witnesses have $n$ bits, giving the
conservative bit bound
\[
 O\bigl(n(S+1)(n+\log(S+1))\bigr)
\]
and space $O((S+1)(n+\log(S+1)))$. Consequently this is a polynomial
algorithm on any subclass with $S\le n^C$ for a fixed constant $C$.
It is only pseudo-polynomial on the full input class.

The promise does not make $S$ polynomially bounded. For $n\ge3$,
take the $n-1$ powers $1,2,\ldots,2^{n-2}$ and append
$2^{n-2}-1$. Their total is $3\cdot2^{n-2}-2<2^n-1$ but grows
exponentially in $n$. They have an easy collision: the last entry
equals the sum of the first $n-2$ entries. This family is not a
hardness proof. It shows that the stated worst-case dynamic-programming
bound cannot be called polynomial merely by invoking the promise.

\paragraph{3. Why a first failure of superincreasing growth is insufficient.}
If an ordered sequence satisfies
$a_i>\sum_{j<i}a_j$ at every step, including $a_1>0$, its prefix
totals obey $S_i\ge2S_{i-1}+1$. Thus $S_n\ge2^n-1$.
The promised input cannot be superincreasing. But the first violation
need not itself identify a collision. The increasing promised input
\[
                         (3,5,6,7,8)
\]
has total $29<31$. Its first violation occurs at six, since
$6\le3+5$, yet the prefix $(3,5,6)$ has the eight distinct subset
sums $0,3,5,6,8,9,11,14$. Overlap of the numeric ranges of two
families of sums does not imply overlap of the discrete families.
The dynamic program handles this gap by retaining all attainable
values, precisely the expensive part that still needs replacement.

\paragraph{4. Modular compression also needs a certificate.}
A collision modulo a small modulus need not be an integer collision.
For example, $(1,2,3,4)$ satisfies the promise, but the singleton
sums one and four agree modulo three and are unequal as integers.
Using a modulus greater than $S$ would make congruence of attainable
sums imply equality, but restoring that many possible residues
reintroduces the same numeric-size dependence. Several small moduli
may filter candidates, yet no bound is proved here ensuring an exact
collision after polynomially many candidate checks.

\paragraph{Verification, first gap and final status.}
The script [S9] runs the witness-producing algorithm on all
promised positive vectors of length at most five with entries at
most six, and on the displayed exponential-total family. It checks
nonempty disjoint witnesses and exact integer equality. The first
remaining gap is reducing the number of stored or examined sums
to polynomial input size without losing a guaranteed collision.
Full promised search problem: UNRESOLVED; the bounded-total subclass
is solved by the explicitly analysed algorithm.
\subsection{Sources}
\begin{itemize}
\item[S1] ULAM.AI and the UnsolvedMath Contributors, supplied problems.json, record 3398 (OPG-163), OpenGarden collection. Catalogue identity and source transcription; CC BY 4.0.
\url{https://huggingface.co/datasets/ulamai/UnsolvedMath}.
\item[S2] Open Problem Garden, Subset-sums equality (pigeonhole version). Original problem and discussion; the mathematical formulation below is expanded from the supplied source record.
\url{http://www.openproblemgarden.org/op/theoretical_computer_science/subset_sums_equality}.
\item[S3] Ce Jin and Hongxun Wu, \emph{A Faster Algorithm for Pigeonhole Equal Sums}, ICALP 2024. \url{https://arxiv.org/abs/2403.19117}.
\item[S4] Deepak Bhati, Pranjal Dutta, Antoine Joux, Mahesh Sreekumar Rajasree and Karol W\k egrzycki, \emph{Pigeonhole Equal Subset Sum: Subexponential Algorithm, Tight Lower Bounds, and Average-Case Analysis}, listed for FOCS 2026 on an author's publication page, accessed 2 October 2026. \url{https://mahe94.github.io/publications/}.
\item[S9] GPT-6 Astra Ultra, exact finite checks accompanying this attempt (2 October 2026). Repository script, Python standard library only. \texttt{scripts/check\_attempt\_batch03\_digraphs\_2026\_10\_02.py}.
\end{itemize}
\end{document}
